\chapter{Methodology}

\section{Dataset}
\paragraph{}
We use a subset of the plant village data set for healthy and diseased tomotoe leaves.

\section{Microservice architecture}
\begin{itemize}
	\item Our image processing pipeline consists of a cluster of docker nodes running image processing services
	\item We use a dockerized ubuntu container with python installation as nodes
	\item Each node runs a flask rest API web service
\end{itemize}

\section{Image processing pipeline}
\begin{itemize}
	\item We use a conventional image classification pipeline
	\item  Preprocessing phase
	\item Segmentation phase
	\item post-procesing phase
	\item feature extraction phase
\end{itemize}

\subsection{Preprocessing Phase}
\begin{itemize}
	\item Color Space Conversion ino HSV (removes luminance issues so more robust to varying lighting conditions)
	\item Filters to make histogram bimodal so foreground and background pixels can be distinguished
	\item apply other transforms like guassian blur, mean histogram filter to make foreground and background pixel intensity levels when changed to gray scale image more differetiated from each other
	\item convert image to gray scale image
\end{itemize}

\subsection{Segmentation Phase}
\begin{itemize}
	\item We apply otsu segmentation thresholding method
	\item due to varying lighting conditions across all images we use adaptive thresholding
\end{itemize}

\subsection{Post-processing Phase}
\begin{itemize}
	\item Artefacts and noise can be introduce from otsu segmentation
	\item We fill in holes, and remove noise and artefacts (apply morphological operations)
	\item Produce a binary mask we can use to isolate the leaf image in the gray level image of leaf
\end{itemize}

\subsection{Feature Extraction Phase}
\begin{itemize}
	\item Textture features are emperically best predictors of tomotoe leaf disease detection
	\item Caculate gray level co-occurence matrix
	\item Entropy, variance etc various other metrics are calculated which make up values in feature vector
\end{itemize}

\section{Classification  Mechanism}
\begin{itemize}
	\item We pass feature vector calculated from gray level co-occurence matrix
	\item This feature vector is fed into SVM classifier which determines whether or not original image leaf has disease or not
	\item This SVM is trained before hand.
	\item Use simple dataset split or 10-fold cross validation
\end{itemize}

\section{Analysis}
\begin{itemize}
	\item Use OpenTelemetry to calculate end-end how much network and compute resources are used
	\item Use OpenTelemetry to calculate how long each image takes to process
	\item Calculate metrics to show reliablity of binary SVM classifier.
		\subitem F measure
		\subitem Accuracy
		\subitem Precesion
		\subitem  true positive rate
		\subitem true negative rate
\end{itemize}